\subsection{Uzasadnienie zastosowania uczenia się ze wzmocnieniem}
\label{sec:uzasadnienie}

Problem nie jest zadaniem klasyfikacji ani jednorazowej optymalizacji, ponieważ:

\begin{enumerate}
  \item \textbf{Decyzje są sekwencyjne} -- misja składa się z kilkudziesięciu
        kroków, a jakość pojedynczego ruchu zależy od całej dalszej trajektorii.
  \item \textbf{Decyzje zmieniają przyszły stan środowiska} -- zużycie energii
        ogranicza dalsze możliwości, a sfotografowanie celu~1 trwale zmienia
        mapę (zamknięcie korytarza $Z$). Kolejność odwiedzania celów ma więc
        znaczenie, a zachłanne wybieranie najbliższego celu może prowadzić
        do sytuacji bez wyjścia.
  \item \textbf{Występuje niepewność} -- wykrycie w strefie radarowej jest
        zdarzeniem losowym, którego prawdopodobieństwo nie jest znane agentowi.
  \item \textbf{Występuje konflikt celów} -- krótsza trasa przez strefę
        radarową oznacza mniejsze zużycie energii, ale większe ryzyko;
        objazd jest bezpieczniejszy, ale może wymagać dodatkowego ładowania.
  \item \textbf{Nagroda jest opóźniona} -- główna nagroda za sukces pojawia
        się dopiero po powrocie do bazy, więc agent musi przypisać zasługę
        wcześniejszym decyzjom (np. odwiedzeniu stacji ładowania).
\end{enumerate}

Te cechy odpowiadają klasycznemu sformułowaniu problemu decyzyjnego
Markowa~\cite{sutton2018}, w którym agent uczy się polityki maksymalizującej
oczekiwany zdyskontowany zwrot bez znajomości modelu przejść.
